%
%****************************************************************************************************
%======================= B E G I N   D O C U M E N T   S E T T I N G S ========================
%****************************************************************************************************
%
\documentclass[a4paper, 12pt, xcolor=dvipsnames]{scrartcl}	% Article als Modul
\usepackage[ngerman]{babel}
\usepackage[onehalfspacing]{setspace}
\usepackage{fancyhdr}
\usepackage[usenames, ps2pdf, dvipsnames,svgnames,table]{xcolor}
\usepackage[utf8x]{inputenc}
\usepackage{wrapfig}   
\usepackage{color}
\usepackage[printonlyused, withpage]{acronym}
\usepackage{ifthen}
\usepackage{lastpage}
\usepackage{currfile}
\usepackage{caption}						% Formatierung der Bildunterschriften
\usepackage{sidecap}
\usepackage{subcaption}
\usepackage{tikz}
\usepackage{calc}
\usepackage{graphicx}
\usepackage{multirow}
\usepackage{longtable}
\usepackage{amsthm}
\usepackage{amsbsy}
\usepackage{amssymb}
\usepackage{array}
\usepackage{chngpage}
\usepackage{tabulary} 					% instead of tabularx
\usepackage{array,longtable,calc,amsmath}
\usepackage{xspace}
\usepackage{pdfpages}
\usepackage{textfit}
\usepackage{tabu}
\usepackage{setspace}
\usepackage{listings}
\pagestyle{fancy}
\usepackage[acronym,toc]{glossaries}
\usepackage{docmute}					% includes parts of files just within the Begin and End Document
\usepackage{blindtext}
\usepackage{lipsum}
%---------------------Switch modulare Printversionen
%
\newboolean{UseWm} 		%Korrektur Wasserzeichen in das PDF drucken
\setboolean{UseWm}{false} 		%Zuweisung eines Wertes                    			% Input file für gemeinsame Package Definitionen
%
%--------------------Set data for document:
\newcommand{\SYear}{2021/22}
\newcommand{\SClass}{5AHET}
\newcommand{\ADatum}{04. April 2022}
\newcommand{\DAName}{E$^2$MILY - Elektrische Hardwaredokumentation}
\newcommand{\DPLNameOnePre}{Vorname1}						% Diplomand 1 als Variablenname, usage: \DPLName...
\newcommand{\DPLNameOneSur}{Nachname1}					% Diplomand 1 als Variablenname, usage: \DPLName...
\newcommand{\DPLNameTwoPre}{Vorname2}						% Diplomand 2 als Variablenname, usage: \DPLName...
\newcommand{\DPLNameTwoSur}{Nachname2}					% Diplomand 2 als Variablenname, usage: \DPLName...
\newcommand{\DPLNameThreePre}{Vorname3}					% Diplomand 3 als Variablenname, usage: \DPLName...
\newcommand{\DPLNameThreeSur}{Nachname3}				% Diplomand 3 als Variablenname, usage: \DPLName...
\newcommand{\DPLNameFourPre}{Vorname4}					% Diplomand 4 als Variablenname, usage: \DPLName...
\newcommand{\DPLNameFourSur}{Nachname4}					% Diplomand 4 als Variablenname, usage: \DPLName...
\newcommand{\BNameOnePre}{Vorname1}							% Betreuer 1 als Variablenname, usage: \DPLName...
\newcommand{\BNameOneSur}{Nachname1}						% Betreuer 1 als Variablenname, usage: \DPLName...
\newcommand{\BNameTwoPre}{Vorname2}						% Betreuer 2 als Variablenname, usage: \DPLName...
\newcommand{\BNameTwoSur}{Nachname2}						% Betreuer 2 als Variablenname, usage: \DPLName...
%
\setboolean{UseWm}{true} 												% Wasserzeichen einschalten
%\newcommand{\WatermarkName}{Korrekturexemplar}		%Text für das Wasserzeichen
%\newcommand{\WatermarkName}{01.01.2021}					%AlternativText für das Wasserzeichen
\newcommand{\WatermarkName}{Template}						%AlternativText für das Wasserzeichen
%
%==============================S E T T I N G S ===================================
%
\definecolor{green}{rgb}{0.8,1.0,0.2}
\definecolor{orange}{rgb}{0.9,0.64,0.25}
\definecolor{kaki}{rgb}{0.74,0.74,0.088}
\definecolor{blue}{rgb}{0.2,0.59,0.94}
\definecolor{dred}{rgb}{0.722,0.18,0.18}
\definecolor{htl}{rgb}{0.15625,0.4375,0.507}
%
%------------------------------Watermark wird mit einer Designvariablen geladen ---------------------------------------------------------------------------------
%
\ifthenelse{\boolean{UseWm}}%
{\usepackage{draftwatermark}\SetWatermarkText{\WatermarkName$~-~$\WatermarkName\xspace$~-~$\WatermarkName} \SetWatermarkLightness{0.95}\SetWatermarkScale{0.5}}{}
%
\graphicspath{{./../03_AUX/01_PICT/}}
\newcommand{\changefont}[3]{\fontfamily{#1} \fontseries{#2} \fontshape{#3} \selectfont} % Change a new font
%
% --------------------Set the figure description:
%
\captionsetup{margin=10pt,font=small, labelfont={color=black!80, bf}, textfont={color=black!60}, format=hang, indention=-1cm}
\captionsetup[wrapfigure]{name=Bild}
\captionsetup[figure]{name=Bild}
% --------------------Set the page margins:
%
\setlength{\voffset}{1.0cm}
\setlength{\textwidth}{17.0cm}
%\addtolength{\textheight}{10.0cm}
\setlength{\textheight}{24.0cm}
\setlength{\topmargin }{-2.3cm}
\setlength{\marginparsep }{0pt}
\setlength{\headsep }{1.2cm}
\setlength{\oddsidemargin}{0.0cm}
\setlength{\evensidemargin}{0.0cm}
\setlength{\marginparwidth}{0pt}
\setlength{\hoffset}{-0.54cm}
\setlength{\headwidth}{17.0cm}
\setlength{\footskip}{30pt}
\setlength{\parindent}{0pt}
%
%

%--------------------Set environment for Header and and Footer:
%
\pagestyle{fancy}
\rhead{\includegraphics[scale=0.06]{/HTL_logo.pdf}}
\chead{}
\lhead{\begin{tabular}{p{12.0cm}} 
\small{\changefont{pag}{b}{sc}\textcolor{black!30}{\DAName}}
\end{tabular} \vspace{0.1cm}}
\lfoot{\small{\changefont{pag}{b}{sc}\textcolor{black!30}{\SYear - \SClass - \DPLNameOneSur, \DPLNameTwoSur, \DPLNameThreeSur, \DPLNameFourSur}}} 
\cfoot{}
\rfoot{\small{\changefont{pag}{b}{sc}\textcolor{black!30}{\thepage\ \textbf{\textcolor{htl!100}I} \pageref{LastPage}}}}
\renewcommand{\headrulewidth}{0.4pt}
\renewcommand{\footrulewidth}{0.4pt}
\let\myHeadrule\headrule
\let\myFootrule\footrule
\renewcommand\headrule{\color{black!30}\myHeadrule }
\renewcommand\footrule{\color{black!30}\myFootrule }
\addto{\captionsngerman}{%
  \renewcommand*{\contentsname}{Inhalt}
  \renewcommand*{\listfigurename}{Abbildungen}
  \renewcommand*{\listtablename}{Abbildungsverzeichnis}
}
%\pagestyle{fancy}
\fancypagestyle{ErsteSeite}{%
   \fancyhf{}%
   \fancyhead[L]{}
}

%\makenoidxglossaries
\renewcommand*\acronymname{Vokabelverzeichnis}
\renewcommand*\glossaryname{Vokabelverzeichnis}

%--------------------Neue Umgebung für M E R K S A T Z:
%
%-------------------- Zählerdefinition:
%
\newcounter{MsNr}
\renewcommand{\theMsNr}{\arabic{MsNr}}
%
\newlength{\boxw}
\newlength{\boxh}
\newlength{\shadowsize}
\newlength{\boxroundness}
\newlength{\tmpa}
\newsavebox{\shadowblockbox}
\setlength{\shadowsize}{6pt}
\setlength{\boxroundness}{3pt}
%
\newenvironment{shadowblock}[1]%
{\begin{lrbox}{\shadowblockbox}\begin{minipage}{#1}}%
{\end{minipage}\end{lrbox}%
\settowidth{\boxw}{\usebox{\shadowblockbox}}%
\settodepth{\tmpa}{\usebox{\shadowblockbox}}%
\settoheight{\boxh}{\usebox{\shadowblockbox}}%
\addtolength{\boxh}{\tmpa}%
\begin{tikzpicture}
\addtolength{\boxw}{\boxroundness * 2}
\addtolength{\boxh}{\boxroundness * 2}
\foreach \x in {0,.05,...,1}
{
\setlength{\tmpa}{\shadowsize * \real{\x}}
\fill[xshift=\shadowsize - 1pt,yshift=-\shadowsize + 
1pt,black,opacity=.04,rounded corners=\boxroundness] 
(\tmpa, \tmpa) rectangle +(\boxw+20 - \tmpa - \tmpa, \boxh+20 - \tmpa - 
\tmpa);
}
\filldraw[fill=yellow!10, draw=kaki!100, rounded corners=\boxroundness] (0, 
0) rectangle (\boxw+20, \boxh+20);
\draw node[xshift=\boxroundness,yshift=\boxroundness,inner sep=10pt,outer 
sep=0pt,anchor=south west] (0,0) {\usebox{\shadowblockbox}};
\end{tikzpicture}}
%
%====================================================================
\newenvironment{Merksatz}
{\begin{shadowblock}{15.8cm}
{\par\vspace*{0.0cm}
\noindent \refstepcounter{MsNr}
\textcolor{kaki!100} {\large{}\textbf{Merke \theMsNr :}} % Hier kann man statt Merke auch etwas anderes hinschreiben
\label{Msl-a}}
}
{\end{shadowblock}}
%
%--------------------End    M E R K S A T Z
% Settings for Codelistings within documentation:
\lstset{ % General setup for the package
	language=C,
	basicstyle=\small\sffamily,
	numbers=left,
 	numberstyle=\tiny,
	frame=tb,
	tabsize=4,
	columns=fixed,
	showstringspaces=false,
	showtabs=false,
	keepspaces,
	commentstyle=\color{red},
	keywordstyle=\color{blue}
}

\setcounter{footnote}{0}

\changefont{pag}{m}{n}                    				% Input file für gemeinsame dokumentenbezogene Definitionen
%
\begin{document}																%ab hier wird der Text im Hauptdokument übernommen
% Verzeichnisse (to be commented out for compilation of main file):
%
%\tableofcontents 																%ACHTUNG: für die Gesamtcompilierung auskommentieren!!
%\newpage 																		%ACHTUNG: für die Gesamtcompilierung auskommentieren!!
%
\changefont{pag}{m}{n}
%
%****************************************************************************************************
%======================== B E G I N  S U B  D O C U M E N T  ==========================
%****************************************************************************************************
%
\section{Regeln für Zitate und Quellenangaben\protect\footnote{Das Quellenverzeichnis ist ein verbindlicher Bestandteil der Diplomarbeit.}}
Ein wesentliches Prinzip wissenschaftlichen Arbeitens ist die Nachvollziehbarkeit der in einer Diplomarbeit
getätigten Aussagen. Werden in einer derartigen schriftlichen Arbeit fremde Quellen verwendet, das heißt zitiert bzw. den eigenen Aussagen zugrunde gelegt, so sind diese Quellen vollständig und korrekt anzugeben.\\

{\begin{tabular}{p{1.0cm}p{12.0cm}}								%Tabelle mit multicolumn
%\rowcolor{orange!10}														%färbige Zeile
	\multicolumn{2}{p{13.0cm}}{\textbf{Derartige Quellen können zum Beispiel sein:}} \\ 
	\textcolor{black!80} {\large\textbf{$\to$}}  & Texte (Bücher, Fachzeitschriften, Firmenunterlagen \dots) \\ 
	\textcolor{black!80} {\large\textbf{$\to$}}  & Filme, Videosequenzen, Radiosendungen \\
	\textcolor{black!80} {\large\textbf{$\to$}}  & Unterrichtsinhalte, Grafiken (Diagramme, Tabellen \dots) \\
	\textcolor{black!80} {\large\textbf{$\to$}}  & Informationen aus dem Internet \\
	\textcolor{black!80} {\large\textbf{$\to$}}  & persönliche Mitteilungen, z.B. externer Fachexperten \\ 
\end{tabular}}

\vspace{0.5cm}																	% vertikaler Abstand
  

\subsection{Zitate}																%Subüberschrift
%
%---------- BEGIN  M e r k s a t z
\begin{Merksatz}
Mit Zitaten belegt man Gedankengänge, Behauptungen und Aussagen.
Sie müssen daher kommentiert und in Beziehung zum konkreten Aspekt der DA gesetzt
werden!
\end{Merksatz}
%---------- END  M e r k s a t z
%

\textbf{Zu vermeiden sind:}
%
\begin{enumerate}
	\linespread {1.0}\changefont{pag}{m}{n}
	\item zu viele und zu umfangreiche Zitate
	\item unnötige Zitate (z. B. technisches Allgemeinwissen)
	\item ungenaue und falsche Zitate
	\item zu wenige Zitate - sind die Ergebnisse wirklich selbst gefunden und geschrieben worden?
	\item aus ihrem Zusammenhang gerissene Zitate
	\linespread {1.25}\changefont{pag}{m}{n}
\end{enumerate}

Das Fehlen korrekter Quellenangaben (z. B. bei Texten, Bilder, Plänen, Zeichnungen, Schaltplänen, Beschreibungen etc.)
kann im Falle der Veröffentlichung der Diplomarbeit schwerwiegende rechtliche Folgen nach sich ziehen und gravierende
finanzielle Auswirkungen (Schadenersatz) bewirken.\\

%
%---------- BEGIN  M e r k s a t z
\begin{Merksatz}
Zitate sind wörtliche Übernahmen aus einem Text und durch Anführungszeichen am Anfang und am Ende als solche zu kennzeichnen!
\end{Merksatz}
%---------- END  M e r k s a t z
%

Es können ganze Sätze oder Satzteile zitiert werden. Zitate können als »wörtliches Zitat« oder als »indirektes Zitat« in den eigenen Text eingefügt werden.

\subsubsection{Das wörtliche Zitat}										%SubSubüberschrift
Das Zitat darf nicht willkürlich aus seinem Textzusammenhang gerissen und sinnentstellend wiedergegeben werden.
Zitate bis zu zwei Zeilen werden in den eigenen Text eingefügt.\\
Zitate über mehr als zwei Zeilen werden eingerückt und im Blocksatz geschrieben. Die Quellenangabe sollte im Anschluss an das Zitat in Klammern angeführt werden. Werden Teile des Textes ausgelassen, so ist das durch Klammern und Auslassungspunkte [...] zu kennzeichnen. Eigene erklärende Anmerkungen im Zitat werden mittels eckiger Klammern [mein Kommentar] markiert. \\

\textbf{Beispiel Langzitat: eingerückt im Blocksatz, Quelle in Kurzform:}\\	%fett gedruckter Text

%---------- BEGIN  LANGZITAT 
\vspace{-0.3cm}																	% vertikaler Abstand

\hspace{1.0cm}																	% horizontaler Abstand
\begin{minipage}{0.935\textwidth}
\linespread {1.25}\changefont{pag}{m}{n}
„Sie wurde zum ersten Mal 1695 in England Wirklichkeit, als das Parlament auf die Zensur [...]
verzichtete. Auf dem Kontinent hat man die Pressefreiheit erst knapp hundert Jahre später
[...] verkündet.“ \cite[S~105]{11_Killinger}
\end{minipage}
\vspace{0.5cm}																	% vertikaler Abstand
%---------- END  T e x t   and  G r a p h i k    in  S u b s e c t i o n
%

\textbf{Beispiel Kurzitat im Text: keine Einrückung, Quelle in Kurzform:}\\
Die Pressefreiheit zählt zu den wichtigsten Kennzeichen einer Demokratie. „Sie wurde zum ersten Mal
1695 in England Wirklichkeit, als das Parlament auf die Zensur [...] verzichtete.“ \cite[S~12]{12_MM}. % Zitat mit Seitenangabe und Kurzzeichen

\subsubsection{Das indirekte Zitat}										%SubSubüberschrift
%---------- BEGIN  M e r k s a t z
\begin{Merksatz}
Der Sinn des Quellentextes darf nicht verändert werden. Indirekte Zitate bleiben ohne Anführungszeichen im Arbeitstext unter Hinzufügung von \cite{BlindzitatKurzzeichen_1}			% Zitat ohne Seitenangabe und Kurzzeichen
\end{Merksatz}
%---------- END  M e r k s a t z
%

\textbf{Beispiel indirektes Zitat:}\\
Die Pressefreiheit zählt zu den wichtigsten Kennzeichen einer Demokratie. Sie wurde in England
1695 zum ersten Mal verkündet. \cite[S~333]{BlindzitatKurzzeichen_1} % Zitat mit Seitenangabe und Kurzzeichen
				
%---------- BEGIN  T e x t   and  G r a p h i k   in  S u b s e c t i o n 
\begin{figure}[h!]
    	\begin{minipage}{0.6\textwidth}
		\begin {center}
		\vspace{0.0cm}
		\centering
		\hspace{-1.0cm}\includegraphics[width=0.8\linewidth]{SecKz_ModKz_Pic_AA}		
		\end {center}
    \end{minipage}%\hfill
    \begin{minipage}{0.4\textwidth}
		\linespread {1.25}\changefont{pag}{m}{n}
		Text im Blindflug zum Mond. Dort kann man sich alles vorstellen.Text im Blindflug zum Mond. Dort kann man sich alles vorstellen.Text im Blindflug zum Mond. Dort kann man sich alles vorstellen.Text im Blindflug zum Mond. Dort kann man sich alles vorstellen.Text im Blindflug zum Mond. Dort kann man sich alles vorstellen.\\
     		\caption{Bildbeschreibung Bildbeschreibung Bildbeschreibung.}
		\label{fig:SecKz_ModKz_Pic_AAb}
    \end{minipage}
\end{figure}
%---------- END  T e x t   and  G r a p h i k    in  S u b s e c t i o n
%----------------------------------------
In Bild \ref{fig:SecKz_ModKz_Pic_AAb} kann man sehen, dass \Blindtext \lipsum[1-1]	%Blindtext zum Auffüllen mit unsinnigen Sätzen
%----------------------------------------
%
\newpage																			%Seitenumbruch am Ende des Subkapitels
%****************************************************************************************************
%============================= E N D   S U B D O C U M E N T  ==========================
%****************************************************************************************************
%
% Verzeichnisse (to be commented out for compilation of main file):
%Im Addendum ist der sogenannte Anhang dargestellt. Beim Verfassen der Dokumentation ist dieser Teil vielleicht noch
%nicht notwendig.
%% Zusammenfassung und Resumee der Diplomarbeit:
\section{Zusammenfassung}

Hier steht die Zusammenfassung und das Resumee der Diplomarbeit!
\newpage

% Anhang zur Diplomarbeit:
\section{Anhang}

Hier stehen zusätzliche Informationen zur Diplomarbeit!
Zum Beispiel können hier Listings oder Teile des Programmes abgebildet sein:

\begin{lstlisting}
// create new document
string path = @"C:\Data\sample.xlsx";
// add title row 
List<string> titleRow = new List<string>();
titleRow.Add("This is the 1st cell");
titleRow.Add("This is the 2nd cell");
openXmlExcel.addTitleRow(titleRow);
// add scatter chart
List<double[]> xValues = new List<double[]>();
List<double[]> yValues = new List<double[]>();
List<string> lineNames = new List<string>();
// first line
xValues.Add(new double[5] { -2, -1, 0, 1, 2 });
yValues.Add(new double[5] { -2, -1, 0, 1, 2 });
lineNames.Add("Line 1");
// second line
xValues.Add(new double[5] { -2, -1, 0, 1, 2 });
yValues.Add(new double[5] { -3, -1, 0, 1, 3 });
lineNames.Add("Line 2");
// save document 
openXmlExcel.saveFile();
\end{lstlisting}

\newpage

\section{Verzeichnisse}
% Bildverzeichnis:
\listoffigures 
%
% Abkürzungsverzeichnis:
%Abkürzungsverzeichnis: Alle Abkürzungen, die bekannt sind finden sich in diesem Verzeichnis.
%Nur diejenigen, die im Dokument auch verwendet werden kommen in die Liste der Compilierung 
%
%der Name der Überschrift kann frei gewählt werden (in diesem Fall: Abkürzungen):
%
%\section*{Acronymverzeichnis}					%Alternative Überschrift für das Verzeichnis
\section*{Abkürzungen}
%
%usage im Dokument: \acs{DA}
%usage im Dokument: \acs{Kurzzeichen}
\begin{acronym}[längstesAcronym]
\setlength{\itemsep}{-0.3 \parsep}
%
% Einträge für Abkürzungen:
%usage im Verzeichnis: \acro{Kurzzeichen}[Inhalt der im Text steht]{Text der als Erklärung im Glossar steht!}
%
\acro{DA}[Diplomarbeit]{Text für die Erklärung der Diplomarbeit}
\acro{ESD}[ESD]{Electrostatic Discharge}
\acro{EMC}[EMC]{Electromagnetic Compatibility}
%
%
% weitere Einträge ...
%
\end{acronym}

%
% Literaturverzeichnis:
%Literaturverzeichnis: Alle Zitate und Veröffentlichungen, die bekannt sind finden sich in diesem Verzeichnis.
%Nur diejenigen, die im Dokument auch verwendet werden kommen in die Liste der Compilierung 
%
%usage im Dokument: \cite{Kurzzeichen}   hier wird im Dokument nur der Text in eckiger Klammer abgebildet
%usage im Dokument: \cite[S~105]{11_Killinger} hier wird auch eine Seitenanzahl dazugestellt
\begin{thebibliography}{22}
%
% Einträge für Literaturverzeichnis:
%usage im Verzeichnis: \bibitem[Werk, Jahrgang oder Auflage]{Kurzzeichen}{Text der als Erklärung im Glossar steht!}
%
   \bibitem[Killinger, 1998]{11_Killinger} Sprache heute : Schriftverkehr; Deutsch für berufsbildende Schulen,\newline 
Seit der Einführung der Bildungsstandards steht der Erwerb fachlicher und sozialer Kompetenzen für Alltag und Beruf im Fokus des Unterrichts.\newline
Autoren: Walter Pirnath, Robert Killinger, Josef Neumüller, 1998
%
\bibitem[vgl. Mustermann, 2020]{12_MM} Skriptum: Test im Zuammenhang mit  ...\newline 
Autoren: Max Mustermann, 2020
%
\bibitem[Werk 1, Jahrgang oder Auflage 1]{BlindzitatKurzzeichen_1}{Text der als Erklärung im Literaturverzeichnis steht!}
%
% weitere Einträge ...
%
\end{thebibliography}
				%ACHTUNG: für die Gesamtcompilierung auskommentieren!!
%Wenn das Addendum auskommentiert ist, dann werden die Zitate im Dokument auch nicht richtig dargestellt.
%Nichts desto Trotz gibt es den Eintrag im Main Dokument auch.
\end{document}


